Soit $\alpha$ un nombre réel non nul tel que $\alpha + \frac{1}{\alpha} \in \mathbb{Z}$.
\begin{enumerate}
  \item Vérifiez que si $\alpha = 2 + \sqrt{3}$, alors $\alpha + \frac{1}{\alpha} = 4$.
  \item Soit $n \geq 1$. On suppose que, pour tout entier $k$ tel que
  $1 \leq k \leq n$, on a $\alpha^k + \frac{1}{\alpha^k} \in \mathbb{Z}$.
  Montrez que cela reste vrai pour $n + 1$, c'est-à-dire que
  $\alpha^{n+1} + \frac{1}{\alpha^{n+1}} \in \mathbb{Z}$.
  (Indication : développer
  $\bigl(\alpha + \frac{1}{\alpha}\bigr)\bigl(\alpha^n + \frac{1}{\alpha^n}\bigr)$ ;
  on rappelle que $\alpha^0 + \frac{1}{\alpha^0} = 2$.)
  \item Concluez (récurrence forte) que pour tout $n \in \mathbb{N}$, $\alpha^n + \frac{1}{\alpha^n} \in \mathbb{Z}$.
\end{enumerate}
